\chapter{Methodologie}

Ce chapitre detaille l'implementation de l'article
\textit{Event-based Real-time Moving Object Detection Based On IMU Ego-motion
Compensation} dans notre environnement ROS 2 \cite{zhao2023imo}.

\section{Motion compensation}
L'etape cle est la compensation de l'ego-mouvement de la camera. Dans
l'article, les auteurs associent des paquets d'evenements a des mesures IMU sur
une fenetre temporelle courte, puis appliquent une fonction de warping non
lineaire issue d'un modele de rotation rigide. L'objectif est de realigner les
evenements du fond statique pour rendre les zones dynamiques plus visibles
\cite{zhao2023imo}.

Toujours selon l'article, la formulation non lineaire ameliore la precision de
compensation (ordre de 10 a 15\%) par rapport a une approximation lineaire,
tout en restant compatible avec une execution temps reel embarquee
\cite{zhao2023imo}.

Sur le plan logiciel, nous sommes partis de la version ROS1 du package de
motion compensation puis nous avons realise une adaptation ROS 2 dans le
package \texttt{datasync\_3\_0}. Le noeud ROS 2:
\begin{itemize}
\item subscribe les evenements et l'IMU;
\item calcule une vitesse angulaire moyenne sur fenetre;
\item produit des evenements compenses, une count-image, une time-image et un
seuil dynamique $\lambda$ publie sur un topic dedie.
\end{itemize}

\section{Objects segmentation}
Apres compensation, le papier exploite le fait que les evenements de fond
stabilises se concentrent spatialement, alors que les evenements d'objets
independants restent plus disperses. La segmentation est faite via une
count-image et une time-image: un pixel est classe dynamique si sa valeur
temporelle normalisee depasse un seuil dynamique $\lambda = a \lVert \omega
\rVert + b$ \cite{zhao2023imo}.

Dans notre implementation ROS 2 (\texttt{project\_ws/src/event\_segmentation}),
le noeud de segmentation:
\begin{itemize}
\item synchronise \texttt{/time\_image} et \texttt{/count\_image};
\item recupere $\lambda$ depuis \texttt{/event\_lambda\_threshold};
\item genere un masque binaire des pixels foreground (avec contraintes de
minimum d'activite et filtrage morphologique optionnel).
\end{itemize}

\section{Objects clustering}
La derniere etape de l'article regroupe les evenements dynamiques en objets.
La logique theorique combine un clustering type DBSCAN (robuste au bruit) avec
une metrique melant position spatio-temporelle 3D et flot optique 2D pour
mieux separer des objets proches dans le plan image \cite{zhao2023imo}.

Dans notre chaine ROS 2 (\texttt{project\_ws/src/event\_clustering}), la
version actuelle realise un regroupement en composantes connexes du masque
foreground, complete par un filtrage temporel de persistance des petites
composantes. Cette adaptation conserve l'esprit de la methode (regroupement
spatio-temporel robuste en temps reel) tout en restant simple a maintenir dans
le pipeline global.

\section{Creation d'une layer Nav2 pour les obstacles dynamiques}
Les tutoriels officiels Nav2 decrivent deux voies complementaires pour injecter
les obstacles dynamiques dans la navigation:
\begin{itemize}
\item creer un plugin de costmap en derivant \texttt{nav2\_costmap\_2d::Layer}
et en implementant les methodes \texttt{onInitialize()},
\texttt{updateBounds()} et \texttt{updateCosts()};
\item exporter ce plugin via \texttt{PLUGINLIB\_EXPORT\_CLASS} et l'activer
dans la liste \texttt{plugins} de la costmap \cite{nav2_plugin_tutorial}.
\end{itemize}

La documentation de configuration Nav2 rappelle aussi que les couches de
costmap peuvent ingerer des sources capteurs via
\texttt{observation\_sources}, notamment en \texttt{LaserScan} ou
\texttt{PointCloud2} \cite{nav2_costmap_config}.

Dans notre projet, l'integration operationnelle est faite par cette seconde
voie: publication des obstacles dynamiques au format \texttt{PointCloud2} sur
\texttt{/dynamic\_obstacles}, puis declaration de cette source dans les
costmaps locale et globale (fichier \texttt{nav2\_diff.yaml}).
